\documentclass[11pt,a4paper]{article}

\usepackage[margin=2.4cm]{geometry}
\usepackage[utf8]{inputenc}
\usepackage[T1]{fontenc}
\usepackage{lmodern}
\usepackage{amsmath,amssymb}
\usepackage{booktabs}
\usepackage{graphicx}
\usepackage{float}

\setlength{\parskip}{0.45em}
\setlength{\parindent}{0pt}

\begin{document}

\noindent
\begin{tabular*}{\textwidth}{@{\extracolsep{\fill}} lr}
\textsc{Universit\`a della Svizzera italiana} & \textsc{Faculty of Economics} \\
\textsc{Teacher: Prof.\ Loriano Mancini} & \textsc{Fall Semester} \\
\textsc{Teaching Assistant: Dardan Gashi} & \\
\end{tabular*}

\vspace{0.6em}
\noindent\textbf{FINANCIAL ECONOMETRICS -- Homework 3} \\
\textbf{Estimating and Testing Multifactor Pricing Models} \\[0.4em]
{\small
Alessia Battistella (25-983-487) \quad Dobri Trifonov (25-983-834) \quad Kyla Kaplan (21-980-321) \\
Tze Qing Ginnie Lee (25-983-719) \quad Lorenzo Comi (25-983-529) \\[0.2em]
19th November 2026
}

\vspace{0.3em}
\hrule
\vspace{1.0em}

\section*{1--2. Data and Descriptive Statistics}

Monthly Fama--French three-factor data (\texttt{Mkt-RF}, \texttt{SMB}, \texttt{HML}, and \texttt{RF}) come from the Kenneth R.\ French Data Library. The CSV contains a monthly block followed by an annual block starting at line 1210; we parse only the monthly section and divide all columns by $100$ to convert from percentages to decimals. Monthly Adjusted Close prices for \texttt{IBM}, \texttt{GE}, and \texttt{JNJ} from January 1980 to December 2019 ($480$ observations) are taken from Yahoo Finance. Because Yahoo's monthly endpoint is missing \texttt{JNJ} data prior to 1985, daily Adjusted Close prices were downloaded and resampled to month-end closing prices, matching the monthly series where both exist.

Computing simple monthly returns $R_{i,t} = P_{i,t}/P_{i,t-1} - 1$ drops January 1980, leaving $T = 479$ observations from February 1980 to December 2019. Slicing the Fama--French factors to the same dates, we form monthly excess returns $y_{i,t} = R_{i,t} - RF_t$. Table~\ref{tab:desc} summarizes the monthly series (in decimals) and Figure~\ref{fig:hist} shows the histograms of monthly excess returns.

\begin{table}[H]
\centering
\caption{Monthly Descriptive Statistics (Feb 1980 -- Dec 2019, $T = 479$, decimals)}
\label{tab:desc}
\small
\begin{tabular}{lrrrrr}
\toprule
Series & Mean & Std.\ Dev. & Min & Median & Max \\
\midrule
IBM ($R_{i,t}$) & 0.0092 & 0.0733 & -0.2619 & 0.0078 & 0.3538 \\
GE ($R_{i,t}$)  & 0.0101 & 0.0730 & -0.2778 & 0.0050 & 0.3421 \\
JNJ ($R_{i,t}$) & 0.0131 & 0.0574 & -0.1734 & 0.0132 & 0.1881 \\
\midrule
IBM ($y_{i,t}$) & 0.0057 & 0.0733 & -0.2647 & 0.0056 & 0.3524 \\
GE ($y_{i,t}$)  & 0.0066 & 0.0728 & -0.2779 & 0.0023 & 0.3400 \\
JNJ ($y_{i,t}$) & 0.0097 & 0.0573 & -0.1812 & 0.0113 & 0.1843 \\
\midrule
Mkt-RF          & 0.0067 & 0.0439 & -0.2319 & 0.0112 & 0.1245 \\
SMB             & 0.0006 & 0.0295 & -0.1739 & 0.0008 & 0.2134 \\
HML             & 0.0024 & 0.0290 & -0.1101 & 0.0009 & 0.1225 \\
RF              & 0.0034 & 0.0029 &  0.0000 & 0.0034 & 0.0135 \\
\bottomrule
\end{tabular}
\end{table}

Among the three stocks, \texttt{JNJ} has the highest average monthly excess return ($0.97\%$) and the lowest volatility ($5.73\%$), while \texttt{IBM} and \texttt{GE} are noticeably more volatile (about $7.3\%$ per month) with wider extreme realizations. The market excess return averages $0.67\%$ per month with a standard deviation of $4.39\%$, whereas \texttt{SMB} and \texttt{HML} have smaller means ($0.06\%$ and $0.24\%$) and lower volatility (about $2.9\%$).

\begin{figure}[H]
    \centering
    \includegraphics[width=\textwidth]{../output/hist_stock_excess_returns.pdf}
    \caption{Histograms of monthly excess returns ($y_{i,t} = R_{i,t} - RF_t$) for IBM, GE, and JNJ.}
    \label{fig:hist}
\end{figure}

\section*{3.1. Estimation}

\subsection*{a. Factor model in vector form}
The Fama--French three-factor model for company $i$ at time $t$ can be written as
\begin{equation*}
    y_{it} = \alpha_i + \underset{1 \times K}{x_t'} \, \underset{K \times 1}{\beta_i} + \varepsilon_{it},
\end{equation*}
where the number of factors is $K = 3$, and the $3 \times 1$ vectors of factors and factor loadings are
\begin{equation*}
    x_t = \begin{bmatrix} RM_t \\ SMB_t \\ HML_t \end{bmatrix}, \qquad
    \beta_i = \begin{bmatrix} \beta_{1,i} \\ \beta_{2,i} \\ \beta_{3,i} \end{bmatrix}.
\end{equation*}

\subsection*{b. Compact matrix form}
Stacking all $T = 479$ monthly observations for firm $i$ gives
\begin{equation*}
    \underset{T \times 1}{y_i} = S\alpha_i + \underset{T \times K}{X} \, \underset{K \times 1}{\beta_i} + \varepsilon_i,
\end{equation*}
where:
\begin{itemize}
    \item $y_i = [y_{i,1}, \dots, y_{i,T}]'$ is the $T \times 1$ ($479 \times 1$) vector of monthly excess returns for firm $i$;
    \item $S = [1, \dots, 1]'$ is a $T \times 1$ vector of ones for the intercept;
    \item $\alpha_i \in \mathbb{R}$ is the scalar intercept (Jensen's alpha) for firm $i$;
    \item $X = [x_1, \dots, x_T]'$ is the $T \times K$ ($479 \times 3$) matrix whose columns are $RM$, $SMB$, and $HML$;
    \item $\beta_i = [\beta_{1,i}, \beta_{2,i}, \beta_{3,i}]'$ is the $K \times 1$ ($3 \times 1$) vector of factor loadings;
    \item $\varepsilon_i = [\varepsilon_{i,1}, \dots, \varepsilon_{i,T}]'$ is the $T \times 1$ vector of i.i.d.\ error terms with $\mathbb{E}[\varepsilon_i] = \mathbf{0}$ and $V(\varepsilon_i) = \sigma_i^2 I_T$.
\end{itemize}

\subsection*{c. Manual OLS estimation}
Defining the augmented design matrix $\tilde{X} = [S,\; X]$ ($479 \times 4$) and parameter vector $\theta_i = [\alpha_i,\; \beta_i']'$ ($4 \times 1$), the model becomes $y_i = \tilde{X}\theta_i + \varepsilon_i$. For each company separately, we compute the OLS estimator manually via matrix algebra (see the accompanying notebook \texttt{main\_analysis.ipynb}):
\begin{equation*}
    \hat{\theta}_i = \begin{bmatrix} \hat{\alpha}_i \\ \hat{\beta}_i \end{bmatrix} = (\tilde{X}'\tilde{X})^{-1}\tilde{X}'y_i.
\end{equation*}
From the residual vector $\hat{\varepsilon}_i = y_i - \tilde{X}\hat{\theta}_i$, the unbiased estimator of the error variance with $T - K - 1 = 475$ degrees of freedom is
\begin{equation*}
    \hat{\sigma}_i^2 = \frac{\hat{\varepsilon}_i'\hat{\varepsilon}_i}{T - K - 1},
\end{equation*}
giving $\hat{V}(\varepsilon_i) = \hat{\sigma}_i^2 I_T$ and parameter covariance matrix $\widehat{\mathrm{Var}}(\hat{\theta}_i) = \hat{\sigma}_i^2(\tilde{X}'\tilde{X})^{-1}$. Standard errors are the square roots of the diagonal entries of $\widehat{\mathrm{Var}}(\hat{\theta}_i)$, and $t$-statistics are evaluated against a $t_{475}$ distribution.

\subsection*{d. Parameter estimates, error covariance, and interpretation}
Table~\ref{tab:ols} reports the estimated intercepts $\hat{\alpha}_i$ and factor loadings $\hat{\beta}_i$ for all three companies, along with standard errors, $t$-statistics, $p$-values, and $R^2$.

\begin{table}[H]
\centering
\caption{Manual OLS Estimates for the Fama--French 3-Factor Model ($T = 479$, $\text{df} = 475$)}
\label{tab:ols}
\small
\begin{tabular}{lccc}
\toprule
 & IBM & GE & JNJ \\
\midrule
$\hat{\alpha}_i$ & $0.0003$ & $-0.0021$ & $0.0058^{***}$ \\
 & $(0.0029)$ & $(0.0024)$ & $(0.0022)$ \\
 & $[t = 0.12,\; p = 0.904]$ & $[t = -0.88,\; p = 0.379]$ & $[t = 2.66,\; p = 0.008]$ \\[0.5em]
$\hat{\beta}_{1,i}$ (\texttt{Mkt-RF}) & $0.8981^{***}$ & $1.2270^{***}$ & $0.6942^{***}$ \\
 & $(0.0671)$ & $(0.0567)$ & $(0.0517)$ \\
 & $[t = 13.39,\; p < 0.001]$ & $[t = 21.62,\; p < 0.001]$ & $[t = 13.44,\; p < 0.001]$ \\[0.5em]
$\hat{\beta}_{2,i}$ (\texttt{SMB}) & $-0.2059^{**}$ & $-0.3057^{***}$ & $-0.5849^{***}$ \\
 & $(0.0993)$ & $(0.0840)$ & $(0.0765)$ \\
 & $[t = -2.07,\; p = 0.039]$ & $[t = -3.64,\; p < 0.001]$ & $[t = -7.64,\; p < 0.001]$ \\[0.5em]
$\hat{\beta}_{3,i}$ (\texttt{HML}) & $-0.2286^{**}$ & $0.2852^{***}$ & $-0.2001^{**}$ \\
 & $(0.1025)$ & $(0.0867)$ & $(0.0789)$ \\
 & $[t = -2.23,\; p = 0.026]$ & $[t = 3.29,\; p = 0.001]$ & $[t = -2.53,\; p = 0.012]$ \\
\midrule
$\hat{\sigma}_i^2$ & 0.003747 & 0.002683 & 0.002223 \\
$\hat{\sigma}_i$   & 0.0612 & 0.0518 & 0.0472 \\
$R^2$              & 0.3071 & 0.4974 & 0.3267 \\
Adjusted $R^2$     & 0.3028 & 0.4943 & 0.3225 \\
\bottomrule
\multicolumn{4}{l}{\footnotesize SE in parentheses; $t$-statistic and $p$-value in brackets. $^{**}p<0.05$, $^{***}p<0.01$.}
\end{tabular}
\end{table}

\textbf{Error variance--covariance estimates $\hat{V}(\varepsilon_i)$:}
Under the single-equation assumption $\varepsilon_i \sim i.i.d.(\mathbf{0}, \sigma_i^2 I_T)$, the estimated $T \times T$ error covariance matrices are
\begin{equation*}
    \hat{V}(\varepsilon_{\text{IBM}}) = 0.003747\, I_{479}, \qquad
    \hat{V}(\varepsilon_{\text{GE}}) = 0.002683\, I_{479}, \qquad
    \hat{V}(\varepsilon_{\text{JNJ}}) = 0.002223\, I_{479}.
\end{equation*}
If we also stack the residuals across the three firms ($\hat{E} = [\hat{\varepsilon}_{\text{IBM}}, \hat{\varepsilon}_{\text{GE}}, \hat{\varepsilon}_{\text{JNJ}}]$), the $3 \times 3$ contemporaneous sample error covariance matrix $\hat{\Sigma}_{\varepsilon} = \frac{1}{475}\hat{E}'\hat{E}$ is:
\begin{equation*}
    \hat{\Sigma}_{\varepsilon} =
    \begin{bmatrix}
        0.003747 & 0.000082 & -0.000128 \\
        0.000082 & 0.002683 &  0.000065 \\
       -0.000128 & 0.000065 &  0.002223
    \end{bmatrix}.
\end{equation*}
The off-diagonal terms are very close to zero (corresponding to residual correlations of $0.026$, $-0.044$, and $0.026$), which shows that the three factors capture almost all common co-movement across these firms.

\textbf{Statistical and economic interpretation:}

\textit{Intercepts ($\hat{\alpha}_i$).} Under the 3-factor model, abnormal returns should be zero ($\alpha_i = 0$). For \texttt{IBM} ($\hat{\alpha} = 0.0003$, $p = 0.904$) and \texttt{GE} ($\hat{\alpha} = -0.0021$, $p = 0.379$), we cannot reject $H_0: \alpha_i = 0$; even though \texttt{GE} underperformed heavily in the latter part of the sample, its average monthly alpha is not statistically different from zero over the full 40-year window. By contrast, \texttt{JNJ} has a positive and statistically significant alpha of $0.0058$ per month (about $7.0\%$ annualized, $t = 2.66$, $p = 0.008$), indicating that \texttt{JNJ} earned returns above what its market, size, and value exposures explain.

\textit{Market exposure ($\hat{\beta}_{1,i}$).} Market betas are large and strongly significant ($p < 0.001$) for all three stocks, and they differ across industries. \texttt{GE} has the highest systematic risk ($\hat{\beta}_1 = 1.2270$), consistent with a cyclical industrial conglomerate that also carried substantial financial-sector exposure through GE Capital. \texttt{IBM} sits just below the market average at $0.8981$. \texttt{JNJ} has a clearly defensive beta of $0.6942$, reflecting stable demand for pharmaceuticals and healthcare products across business cycles.

\textit{Size exposure ($\hat{\beta}_{2,i}$).} All three companies load negatively and significantly on \texttt{SMB}: $-0.2059$ for \texttt{IBM} ($p = 0.039$), $-0.3057$ for \texttt{GE} ($p < 0.001$), and $-0.5849$ for \texttt{JNJ} ($p < 0.001$). Because \texttt{SMB} is long small firms and short large firms, negative loadings are expected for Dow Jones mega-cap companies. The exposure is especially large in magnitude for \texttt{JNJ} ($-0.58$), reflecting its steady large-cap profile over the sample.

\textit{Value--growth exposure ($\hat{\beta}_{3,i}$).} The \texttt{HML} loadings are smaller in magnitude than the market betas, statistically significant at the $5\%$ level for all three firms, and differ in sign. \texttt{GE} loads positively on \texttt{HML} ($+0.2852$, $p = 0.001$), behaving like a value stock due to its capital-intensive industrial and financial assets. \texttt{IBM} ($-0.2286$, $p = 0.026$) and \texttt{JNJ} ($-0.2001$, $p = 0.012$) both load negatively, tilting modestly toward growth stocks whose valuations rely more heavily on R\&D, patents, and brand intangibles rather than book assets.

\textit{Fit and idiosyncratic risk.} The three factors explain nearly half of the return variance for \texttt{GE} ($R^2 = 0.4974$) and roughly one third for \texttt{JNJ} ($0.3267$) and \texttt{IBM} ($0.3071$). Firm-specific monthly residual volatility is largest for \texttt{IBM} ($\hat{\sigma} = 6.12\%$), reflecting major company-specific shifts over the period, followed by \texttt{GE} ($5.18\%$) and \texttt{JNJ} ($4.72\%$).

\end{document}
